
\subsubsection{XXX 结果的分析}

% 本节对求解结果进行深入分析。
% 关键内容:
%   1. 关键发现(主方法 vs 对比方法)
%   2. 与领域直觉/理论的一致性
%   3. 局限与启示
%   4. 总结桥接到下一节

研究结果表明,【关键结论】。这一发现与【领域直觉/理论】一致,【进一步解释】。

【进一步说明】。然而,【简单方法】对【响应变量】变异的解释力有限(【指标】较低),提示实际关系中存在复杂的非线性成分和未观测的混杂因素。

通过【主方法】与【对比方法】的对比,可以看出:

\begin{itemize}[itemindent=2em]
  \item \textbf{【因素 1】:}【说明其对响应变量的影响方向和大小】。
  \item \textbf{【因素 2】:}【说明其对响应变量的影响方向和大小】。
  \item \textbf{【因素 3】:}【说明其对响应变量的影响方向和大小】。
\end{itemize}

% 关键结果对比表(可省略,已在 5.1.3 中给出)
% \begin{longtable}{...}
%   ...
% \end{longtable}

综上所述,本问建立了【响应变量】与【自变量】的【模型类型】关系模型,识别出【主要因素】是主要影响因素。模型【整体显著 / 解释力有限】,【需要更复杂的建模方法以提高拟合优度 / 已达到预期精度】。
